\documentclass[10pt,a4paper]{article}

\usepackage[utf8]{inputenc}
\usepackage[T1]{fontenc}
\usepackage{lmodern}
\usepackage[spanish]{babel}
\usepackage[a4paper,top=1.0cm,bottom=1.0cm,left=1.55cm,right=1.55cm]{geometry}
\usepackage{enumitem}
\usepackage{titlesec}
\usepackage{xcolor}
\usepackage[hidelinks]{hyperref}
\usepackage{microtype}

\pagestyle{empty}
\setlength{\parindent}{0pt}
\setlength{\parskip}{1.5pt}
\setlength{\emergencystretch}{2em}
\linespread{1.0}
\color{black}

\titleformat{\section}{\large\bfseries\color{black}}{}{0pt}{}[\titlerule]
\titlespacing*{\section}{0pt}{4pt}{2pt}
\setlist[itemize]{leftmargin=1.2em,itemsep=1pt,topsep=1pt,parsep=0pt,partopsep=0pt}

\begin{document}

\begin{center}
{\LARGE\bfseries ABDIEL CAMARGO}\\[2pt]
{\normalsize\bfseries CIBERSEGURIDAD \;|\; CLOUD \;|\; INGENIERÍA DE IA}\\[3pt]
{\small David, Panamá \;|\; (+507) 6786-2691 \;|\; \href{mailto:camargoabdiel15@gmail.com}{camargoabdiel15@gmail.com}}\\[1pt]
{\small Portafolio: \href{https://abdiellzz.github.io/portafolio/}{abdiellzz.github.io/portafolio} \;|\; \href{https://www.linkedin.com/in/camargoabdiel/}{linkedin.com/in/camargoabdiel}}
\end{center}

\vspace{-4pt}

\section*{Perfil Profesional}
Estudiante de Ingeniería en Ciberseguridad que combina base técnica (AWS, IAM, Dynatrace, Python, IA local con Ollama) con trato directo con usuarios: doy soporte técnico, coordino equipos y presento ideas en público. Recibí un reconocimiento en la Presidencia de la República de Panamá de manos de la Primera Dama, Maricel Cohen de Mulino, por el Programa de IA, Innovación y Agilidad (MIDES y EUDE). 4.º lugar en el CTF Nacional de Ciberseguridad y \textbf{1.er lugar} en la Feria del MIDES con ALLEN. Busco aportar en soporte de TI, cloud o ciberseguridad.

\section*{Habilidades Técnicas}
\textbf{Sistemas y Soporte:} Windows, fundamentos de Linux, Microsoft 365, mantenimiento de hardware, configuración y soporte técnico.\\
\textbf{Redes y Observabilidad:} Fundamentos de redes, Dynatrace, DQL, análisis de logs y monitoreo de servicios.\\
\textbf{Programación, Datos e IA:} Python, SQL, Git/GitHub, automatización, Ollama y LLM locales.\\
\textbf{Nube e Infraestructura:} AWS (servicios principales, IAM), Azure, Kubernetes y Terraform.

\section*{Experiencia Profesional}
\textbf{Asistente de Soporte de TI \;|\; Universidad del Istmo} \hfill \textit{Mayo 2026 -- Presente}
\begin{itemize}
  \item Diagnostiqué y resolví problemas de hardware, software y conectividad.
  \item Mantuve, actualicé y configuré entornos Windows y herramientas de seguridad.
  \item Desarrollé herramientas internas para inventario de activos de TI, seguimiento de actualizaciones y monitoreo de sistemas, en coordinación con equipos multidisciplinarios.
\end{itemize}

\textbf{Gerente de Auditoría \;|\; AIESEC} \hfill \textit{Ago. 2026 -- Ene. 2027}
\begin{itemize}
  \item Gestioné la documentación requerida por la entidad auditora para el programa de intercambio internacional, manteniendo los expedientes de cada pasante listos para auditoría.
  \item Recopilé y verifiqué evidencias de la implementación de los estándares de intercambio, demostrando el cumplimiento del programa.
  \item Mantuve plataformas de seguimiento actualizadas, resguardando la experiencia de los pasantes y el cumplimiento del marco migratorio y legal.
\end{itemize}

\textbf{Pasante de TI \;|\; Institute College Panamá / CSCYT} \hfill \textit{Oct. -- Dic. 2025}
\begin{itemize}
  \item Brindé soporte de redes y asistencia técnica; administré bases de datos y sistemas de información, gestionando documentación con Microsoft Office y procesos administrativos.
\end{itemize}

\section*{Proyectos Técnicos}
\begin{itemize}
  \item \textbf{ALLEN \;|\; Desarrollador, 2026.} Construí y probé con usuarios un prototipo de herramienta educativa local con IA (Python, Ollama, SQLite, JavaScript). \textbf{1.er lugar} en el pitch de la Feria del MIDES (Boquete, oct. 2026).
  \item \textbf{7.ª Hackathon Copa Airlines \;|\; Microsoft y Dynatrace, 8--9 oct. 2026.} Dynatrace/DQL para logs, análisis de tráfico y monitoreo de servicios en Kubernetes; soluciones desplegadas en Azure y Terraform.
  \item \textbf{Laboratorio Introducción a AWS IAM \;|\; Portafolio.} Configuré usuarios, grupos y políticas y probé accesos aplicando mínimo privilegio; documentado con video guía y capturas.
\end{itemize}

\section*{Educación}
\textbf{Licenciatura en Ingeniería en Ciberseguridad} \hfill \textit{En curso, 3.er trimestre}\\
Universidad del Istmo\\[2pt]
\textbf{Programa de Inteligencia Artificial, Innovación y Agilidad para la Transformación Profesional}\\
EUDE European Business School \;|\; 20 horas \;|\; 21 oct.--13 nov. 2026, completado. Diploma entregado en la Presidencia de la República de Panamá.

\section*{Certificaciones y Reconocimientos}
{\small
\textbf{4.º Lugar Nacional — CTF de Ciberseguridad} (8,385 pts) — Min. de Seguridad Pública / ITSE (23 oct. 2026)\\
\textbf{1.er Lugar — Feria de Innovación y Emprendimiento del MIDES, Boquete} (ALLEN) — 3 oct. 2026\\
\textbf{Finalista — hackIAthon de Viamatica ONDU} (equipo JOHNPORK) — 16 oct. 2026\\
\textbf{AWS Academy Cloud Foundations} (en curso, oct. 2026) \;|\; \textbf{IA Generativa con AWS} (ago. 2026) — AWS\\
\textbf{Introducción a la infraestructura en la nube} (ago. 2026) \;|\; \textbf{Administrador de Microsoft 365} (feb. 2026) — Microsoft\\
\textbf{Linux Unhatched} — Cisco Networking Academy / UTP (sep. 2026)\\
\textbf{Fundamentos de la ciberseguridad} — Google / Coursera (ene. 2026)
}

\section*{Idiomas y Habilidades Profesionales}
{\small\textbf{Idiomas:} Español (nativo) \;|\; Inglés (B1 intermedio, EF SET 44/100 — nov. 2026).\\
Comunicación efectiva, trabajo en equipo, empatía y orientación al usuario, liderazgo, proactividad, responsabilidad, organización y gestión del tiempo, atención al detalle, pensamiento analítico, resolución de problemas, adaptabilidad y aprendizaje continuo.}

\end{document}
